\documentclass{article}
\usepackage{amsmath,amssymb,amsthm}
\usepackage{algorithm}
\usepackage{algpseudocode}
\usepackage{hyperref}
\usepackage{xcolor}
\newtheorem{theorem}{Theorem}
\newtheorem{lemma}{Lemma}
\newcommand{\E}{\mathbb{E}}
\newcommand{\R}{\mathbb{R}}
\newcommand{\inner}[2]{\langle #1,#2\rangle}
\newcommand{\argmin}{\operatorname*{arg\,min}}
\newcommand{\D}[2]{D_{h}(#1,#2)}
\newcommand{\grad}{\nabla}
\newcommand{\prox}{\operatorname{prox}}
\begin{document}

\section*{Algorithm Name}
\textbf{Mirror Descent with Bregman Proximal Update} (also known as the Bregman gradient method).

-----------------------------------------------------------------------

\section*{Mathematical Formulation}
Let $f:\R^{n}\rightarrow\R$ be a convex differentiable function.  
Let $h:\R^{n}\rightarrow\R$ be a \emph{prox‑function}, i.e., a continuously differentiable, $\sigma$–strongly convex function w.r.t. a norm $\|\cdot\|$:
\[
h(y)\ge h(x)+\inner{\grad h(x)}{y-x}+\frac{\sigma}{2}\|y-x\|^{2},\qquad\forall x,y\in\R^{n}.
\]
The associated Bregman divergence is
\[
\D{y}{x}=h(y)-h(x)-\inner{\grad h(x)}{y-x}.
\]

Given a stepsize (learning rate) sequence $\{\eta_k\}_{k\ge0}$ with $\eta_k>0$, the update at iteration $k$ is
\begin{equation}
\boxed{
x_{k+1}= \argmin_{x\in\R^{n}}\Big\{\inner{\grad f(x_k)}{x}+ \frac{1}{\eta_k}\,\D{x}{x_k}\Big\}.
}
\label{eq:update}
\end{equation}
When $h(x)=\frac12\|x\|^{2}$, $D_h$ reduces to the Euclidean squared distance and \eqref{eq:update} becomes the classical gradient step.

-----------------------------------------------------------------------

\section*{Pseudocode}
\begin{algorithm}[H]
\caption{Mirror Descent with Bregman Proximal Update}
\begin{algorithmic}[1]
\State \textbf{Input:} initial point $x_0\in\R^{n}$, stepsize schedule $\{\eta_k\}_{k\ge0}$, prox‑function $h$.
\For{$k=0,1,2,\dots$}
    \State Compute gradient $g_k=\grad f(x_k)$.
    \State Compute the next iterate
    \[
    x_{k+1}= \argmin_{x}\Big\{\inner{g_k}{x}+ \frac{1}{\eta_k}\big( h(x)-h(x_k)-\inner{\grad h(x_k)}{x-x_k}\big)\Big\}.
    \]
\EndFor
\State \textbf{Output:} $x_T$ (or the average $\bar x_T=\frac1T\sum_{k=0}^{T-1}x_k$).
\end{algorithmic}
\end{algorithm}

-----------------------------------------------------------------------

\section*{Dry Run Example}
Consider the one‑dimensional problem 
\[
f(w)=(w-3)^{2},\qquad \grad f(w)=2(w-3).
\]
Choose the Euclidean prox‑function $h(w)=\tfrac12 w^{2}$ (hence $D_h(u,v)=\tfrac12 (u-v)^{2}$) and a constant stepsize $\eta_k=\eta=0.1$.  
Initialize $w_{0}=0$.

\medskip
\textbf{Iteration 0}
\begin{align*}
g_{0}&=\grad f(w_{0})=2(0-3)=-6,\\
w_{1}&=\argmin_{w}\Big\{-6w+\frac{1}{0.1}\,\tfrac12 (w-0)^{2}\Big\}
        =\argmin_{w}\Big\{-6w+5w^{2}\Big\}\\
      &\stackrel{\text{first‑order}}{=}\frac{6}{10}=0.6,\\
f(w_{1})&=(0.6-3)^{2}=5.76.
\end{align*}

\textbf{Iteration 1}
\begin{align*}
g_{1}&=2(w_{1}-3)=2(0.6-3)=-4.8,\\
w_{2}&=\argmin_{w}\Big\{-4.8w+5(w-0.6)^{2}\Big\}
        =\frac{4.8+5\cdot0.6}{10}=0.96,\\
f(w_{2})&=(0.96-3)^{2}=4.1616.
\end{align*}

\textbf{Iteration 2}
\begin{align*}
g_{2}&=2(w_{2}-3)=2(0.96-3)=-4.08,\\
w_{3}&=\argmin_{w}\Big\{-4.08w+5(w-0.96)^{2}\Big\}
        =\frac{4.08+5\cdot0.96}{10}=1.248,\\
f(w_{3})&=(1.248-3)^{2}=3.066.
\end{align*}

The objective value decreases monotonically, illustrating the descent property of the algorithm.

-----------------------------------------------------------------------

\section*{Assumptions}
\begin{enumerate}
\item[(A1)] \textbf{Convexity of $f$.}  
      $f$ is convex on $\R^{n}$:
      \[
      f(y)\ge f(x)+\inner{\grad f(x)}{y-x},\qquad\forall x,y.
      \]

\item[(A2)] \textbf{Smoothness w.r.t. $D_h$.}  
      There exists $L>0$ such that for all $x,y$,
      \[
      f(y)\le f(x)+\inner{\grad f(x)}{y-x}+L\,\D{y}{x}. \tag{1}
      \]

\item[(A3)] \textbf{Strong convexity of $h$.}  
      $h$ is $\sigma$‑strongly convex w.r.t. $\|\cdot\|$:
      \[
      \D{y}{x}\ge\frac{\sigma}{2}\|y-x\|^{2},\qquad\forall x,y. \tag{2}
      \]

\item[(A4)] \textbf{Stepsize.}  
      The stepsizes satisfy $0<\eta_k\le \frac{1}{L}$ for all $k$ (constant stepsize $\eta\le 1/L$ is a special case).
\end{enumerate}

-----------------------------------------------------------------------

\section*{Main Convergence Theorem}
\begin{theorem}[Convergence of Mirror Descent]\label{thm:conv}
Assume (A1)–(A4) and let $x^{\star}\in\arg\min f$.  
Define the Bregman diameter of the level set $\mathcal{X}:=\{x\mid f(x)\le f(x_{0})\}$ as
\[
\Omega:=\sup_{x\in\mathcal{X}}\D{x}{x_{0}}<\infty.
\]
Then for any $T\ge1$ the iterates generated by \eqref{eq:update} satisfy
\[
f(\bar x_T)-f(x^{\star})\;\le\;\frac{\Omega}{\eta T},
\qquad\text{where }\;\bar x_T:=\frac1T\sum_{k=0}^{T-1}x_k.
\tag{3}
\]
Consequently, with a constant stepsize $\eta=1/L$ we obtain the $O(1/T)$ rate.

If, in addition, $f$ is $\mu$‑strongly convex (i.e. $f(y)\ge f(x)+\inner{\grad f(x)}{y-x}+\frac{\mu}{2}\|y-x\|^{2}$), then
\[
\|x_{k}-x^{\star}\|^{2}\;\le\;\Big(1-\eta\mu\sigma\Big)^{k}\,\frac{2\Omega}{\sigma},
\tag{4}
\]
i.e. a linear (geometric) convergence.
\end{theorem}

-----------------------------------------------------------------------

\section*{Theoretical Properties}
\begin{itemize}
\item \textbf{Stability.} The update \eqref{eq:update} is a proximal point in the dual space; the strong convexity of $h$ guarantees existence and uniqueness of $x_{k+1}$ and prevents oscillations.
\item \textbf{Hyperparameter Sensitivity.} The bound \eqref{3} depends on the product $\eta T$; choosing $\eta$ close to $1/L$ maximizes the progress per iteration while preserving the feasibility of the smoothness inequality (1). Over‑large $\eta$ may violate (1) and break the guarantee.
\item \textbf{Comparison to Euclidean Gradient Descent.} When $h(x)=\frac12\|x\|^{2}$, $D_h$ reduces to $\frac12\|x-y\|^{2}$ and Theorem~\ref{thm:conv} recovers the classical $O(1/T)$ rate for convex smooth functions and the $O((1-\eta\mu)^{k})$ linear rate for strongly convex functions.
\item \textbf{Adaptivity.} By selecting a geometry‑adapted $h$ (e.g. entropy for the simplex), the Bregman diameter $\Omega$ can be dramatically smaller than a Euclidean diameter, leading to tighter bounds in practice.
\end{itemize}

-----------------------------------------------------------------------

\section*{Discussion}
The theorem shows that Mirror Descent attains the optimal $O(1/T)$ rate for convex smooth optimization under the Bregman‑smoothness condition (A2). The proof relies only on first‑order information and the three‑point identity of Bregman divergences; no stochasticity or variance reduction is required. The linear rate \eqref{4} demonstrates that strong convexity of $f$ together with the strong convexity of the prox‑function yields a contraction factor $1-\eta\mu\sigma$, mirroring the classical analysis but expressed in the geometry induced by $h$.

-----------------------------------------------------------------------

\section*{Preliminaries (Definitions and Lemmas)}
\begin{lemma}[Three‑Point Identity]\label{lem:3pt}
For any $x,y,z\in\R^{n}$,
\[
\inner{\grad h(z)-\grad h(y)}{x-y}= \D{x}{y}+ \D{y}{z}-\D{x}{z}. \tag{5}
\]
\end{lemma}
\begin{proof}
Direct expansion of the definition of $D_h$ yields \eqref{5}. \hfill\qedsymbol
\end{proof}

\begin{lemma}[Bregman Strong Convexity]\label{lem:strong}
If $h$ is $\sigma$‑strongly convex then for all $x,y$,
\[
\D{x}{y}\ge\frac{\sigma}{2}\|x-y\|^{2}. \tag{6}
\]
\end{lemma}
\begin{proof}
Immediate from the definition of strong convexity. \hfill\qedsymbol
\end{proof}

\begin{lemma}[Smoothness w.r.t. $D_h$]\label{lem:smooth}
Assumption (A2) is equivalent to
\[
\inner{\grad f(x)-\grad f(y)}{x-y}\le L\,\D{x}{y},\qquad\forall x,y. \tag{7}
\]
\end{lemma}
\begin{proof}
Subtract the two inequalities in (1) written for $(x,y)$ and $(y,x)$ and rearrange. \hfill\qedsymbol
\end{proof}

-----------------------------------------------------------------------

\section*{Main Proof (Step‑by‑Step Derivation)}
We prove the bound \eqref{3}.  
Let $x^{\star}\in\arg\min f$ be a minimizer. For each iteration $k$ we start from the optimality condition of \eqref{eq:update}.

\medskip
\noindent\textbf{[Step 1]} (Optimality condition).  
Since $x_{k+1}$ minimizes a convex function, the first‑order condition gives
\[
\inner{\grad f(x_k)}{x-x_{k+1}}+\frac{1}{\eta_k}\inner{\grad h(x_{k+1})-\grad h(x_k)}{x-x_{k+1}}\ge0,
\qquad\forall x\in\R^{n}. \tag{8}
\]

\noindent\textbf{[Step 2]} (Choose $x=x^{\star}$).  
Insert $x^{\star}$ into \eqref{8}:
\[
\inner{\grad f(x_k)}{x^{\star}-x_{k+1}}+\frac{1}{\eta_k}\inner{\grad h(x_{k+1})-\grad h(x_k)}{x^{\star}-x_{k+1}}\ge0. \tag{9}
\]

\noindent\textbf{[Step 3]} (Rearrange the second inner product using Lemma~\ref{lem:3pt}).  
Apply identity \eqref{5} with $(x,y,z)=(x^{\star},x_{k+1},x_k)$:
\[
\inner{\grad h(x_{k+1})-\grad h(x_k)}{x^{\star}-x_{k+1}}
= \D{x^{\star}}{x_k}-\D{x^{\star}}{x_{k+1}}-\D{x_{k+1}}{x_k}. \tag{10}
\]

\noindent\textbf{[Step 4]} (Plug \eqref{10} into \eqref{9}).  
We obtain
\[
\inner{\grad f(x_k)}{x^{\star}-x_{k+1}}
+\frac{1}{\eta_k}\Big(\D{x^{\star}}{x_k}-\D{x^{\star}}{x_{k+1}}-\D{x_{k+1}}{x_k}\Big)\ge0.
\tag{11}
\]

\noindent\textbf{[Step 5]} (Use convexity of $f$).  
Convexity (A1) yields
\[
f(x_k)-f(x^{\star})\le\inner{\grad f(x_k)}{x_k-x^{\star}}.
\tag{12}
\]
Rewriting $\inner{\grad f(x_k)}{x^{\star}-x_{k+1}}
= \inner{\grad f(x_k)}{x^{\star}-x_k}+\inner{\grad f(x_k)}{x_k-x_{k+1}}$ and applying \eqref{12} gives
\[
\inner{\grad f(x_k)}{x^{\star}-x_{k+1}}
\le f(x^{\star})-f(x_k)+\inner{\grad f(x_k)}{x_k-x_{k+1}}.
\tag{13}
\]

\noindent\textbf{[Step 6]} (Upper bound the term $\inner{\grad f(x_k)}{x_k-x_{k+1}}$ using smoothness).  
From Lemma~\ref{lem:smooth} (inequality \eqref{7}) with $(x,y)=(x_k,x_{k+1})$ we have
\[
\inner{\grad f(x_k)-\grad f(x_{k+1})}{x_k-x_{k+1}}\le L\,\D{x_k}{x_{k+1}}.
\tag{14}
\]
Rearranging,
\[
\inner{\grad f(x_k)}{x_k-x_{k+1}}\le \inner{\grad f(x_{k+1})}{x_k-x_{k+1}}+L\,\D{x_k}{x_{k+1}}.
\tag{15}
\]
Since $x_{k+1}$ is the minimizer of \eqref{eq:update}, the first‑order optimality condition with $x=x_k$ gives
\[
\inner{\grad f(x_k)}{x_{k+1}-x_k}+\frac{1}{\eta_k}\inner{\grad h(x_{k+1})-\grad h(x_k)}{x_{k+1}-x_k}\le0,
\]
which after multiplying by $-1$ reads
\[
\inner{\grad f(x_k)}{x_k-x_{k+1}}\le\frac{1}{\eta_k}\inner{\grad h(x_k)-\grad h(x_{k+1})}{x_k-x_{k+1}}.
\tag{16}
\]
Using Lemma~\ref{lem:3pt} with $(x,y,z)=(x_k,x_{k+1},x_k)$ we obtain
\[
\inner{\grad h(x_k)-\grad h(x_{k+1})}{x_k-x_{k+1}} = \D{x_k}{x_{k+1}}+\D{x_{k+1}}{x_k}
= 2\,\D{x_{k+1}}{x_k},   \tag{17}
\]
where the symmetry $\D{x_k}{x_{k+1}}=\D{x_{k+1}}{x_k}$ follows from the definition of $D_h$ when $h$ is twice differentiable.  
Insert \eqref{17} into \eqref{16}:
\[
\inner{\grad f(x_k)}{x_k-x_{k+1}}\le\frac{2}{\eta_k}\,\D{x_{k+1}}{x_k}. \tag{18}
\]

\noindent\textbf{[Step 7]} (Combine \eqref{13}, \eqref{18} and \eqref{11}).  
From \eqref{13} and \eqref{18} we have
\[
\inner{\grad f(x_k)}{x^{\star}-x_{k+1}}
\le f(x^{\star})-f(x_k)+\frac{2}{\eta_k}\,\D{x_{k+1}}{x_k}. \tag{19}
\]
Plugging \eqref{19} into the left‑hand side of \eqref{11} yields
\[
f(x^{\star})-f(x_k)+\frac{2}{\eta_k}\,\D{x_{k+1}}{x_k}
+\frac{1}{\eta_k}\big(\D{x^{\star}}{x_k}-\D{x^{\star}}{x_{k+1}}-\D{x_{k+1}}{x_k}\big)\ge0.
\tag{20}
\]

\noindent\textbf{[Step 8]} (Simplify \eqref{20}).  
Collecting the Bregman terms:
\[
f(x_k)-f(x^{\star})\le\frac{1}{\eta_k}\big(\D{x^{\star}}{x_k}-\D{x^{\star}}{x_{k+1}}\big)
+\Big(\frac{1}{\eta_k}-\frac{2}{\eta_k}\Big)\D{x_{k+1}}{x_k).
\tag{21}
\]
Since $\eta_k\le 1/L$, the coefficient $\frac{1}{\eta_k}-\frac{2}{\eta_k}= -\frac{1}{\eta_k}$ is non‑positive; discarding the non‑positive term gives the inequality
\[
f(x_k)-f(x^{\star})\le\frac{1}{\eta_k}\big(\D{x^{\star}}{x_k}-\D{x^{\star}}{x_{k+1}}\big). \tag{22}
\]

\noindent\textbf{[Step 9]} (Telescoping sum).  
Sum \eqref{22} over $k=0,\dots,T-1$:
\[
\sum_{k=0}^{T-1}\big(f(x_k)-f(x^{\star})\big)
\le\frac{1}{\eta}\big(\D{x^{\star}}{x_0}-\D{x^{\star}}{x_T}\big)
\le\frac{\Omega}{\eta}, \tag{23}
\]
where we used the constant stepsize $\eta_k\equiv\eta$ and the definition of $\Omega$.

\noindent\textbf{[Step 10]} (Convexity of the average).  
By convexity of $f$,
\[
f(\bar x_T)-f(x^{\star})\le\frac{1}{T}\sum_{k=0}^{T-1}\big(f(x_k)-f(x^{\star})\big). \tag{24}
\]

\noindent\textbf{[Step 11]} (Combine \eqref{23} and \eqref{24}).  
We obtain the desired bound
\[
f(\bar x_T)-f(x^{\star})\le\frac{\Omega}{\eta T}. \tag{25}
\]
This completes the proof of the $O(1/T)$ rate.

\medskip
\noindent\textbf{Proof of the linear rate \eqref{4}.}  
Assume additionally that $f$ is $\mu$‑strongly convex:
\[
f(y)\ge f(x)+\inner{\grad f(x)}{y-x}+\frac{\mu}{2}\|y-x\|^{2}. \tag{26}
\]
Apply the three‑point identity \eqref{5} with $x=x^{\star}$, $y=x_{k+1}$, $z=x_k$ and use (26) together with Lemma~\ref{lem:strong} (6):
\[
\begin{aligned}
\D{x^{\star}}{x_{k+1}}
&\le \D{x^{\star}}{x_k}-\eta\big(f(x_k)-f(x^{\star})\big)-\eta\frac{\mu}{2}\|x_k-x^{\star}\|^{2}\\
&\le \D{x^{\star}}{x_k}-\eta\frac{\mu\sigma}{2}\|x_k-x^{\star}\|^{2} \qquad\text{(by }\eqref{6}\text{)}\\
&\le \big(1-\eta\mu\sigma\big)\,\D{x^{\star}}{x_k}.
\end{aligned}
\tag{27}
\]
Iterating (27) yields
\[
\D{x^{\star}}{x_{k}}\le\big(1-\eta\mu\sigma\big)^{k}\D{x^{\star}}{x_0}
\le\big(1-\eta\mu\sigma\big)^{k}\,\Omega. \tag{28}
\]
Finally, by Lemma~\ref{lem:strong},
\[
\|x_k-x^{\star}\|^{2}\le\frac{2}{\sigma}\D{x^{\star}}{x_k}
\le\frac{2\Omega}{\sigma}\big(1-\eta\mu\sigma\big)^{k},
\]
which is exactly \eqref{4}. \hfill\qedsymbol

-----------------------------------------------------------------------

\section*{Rate Derivation (With Constants)}
For a constant stepsize $\eta=1/L$ the bound \eqref{25} becomes
\[
f(\bar x_T)-f(x^{\star})\le\frac{L\,\Omega}{T}.
\]
If the level‑set diameter $\Omega$ is known (e.g. $\Omega\le \frac{R^{2}}{2}$ for Euclidean $h$ on a ball of radius $R$), the constant factor can be made explicit.  
In the strongly convex case the contraction factor is $q:=1-\frac{\mu\sigma}{L}<1$, giving
\[
\|x_k-x^{\star}\|^{2}\le\frac{2\Omega}{\sigma}\,q^{k}.
\]

-----------------------------------------------------------------------

\section*{Special Cases}
\begin{itemize}
\item \textbf{Euclidean setup.} $h(x)=\frac12\|x\|^{2}$, $\sigma=1$, $D_h(u,v)=\frac12\|u-v\|^{2}$. The theorem recovers the classic gradient descent rates.
\item \textbf{Entropy prox‑function on the simplex.} $h(x)=\sum_{i}x_i\log x_i$, $\sigma=1$ w.r.t. $\ell_{1}$ norm, $D_h$ is the Kullback–Leibler divergence. The bound involves the KL‑diameter of the feasible set, which can be far smaller than a Euclidean diameter.
\item \textbf{Non‑constant stepsizes.} If $\eta_k=\frac{c}{k+1}$ with $c\le \frac{1}{L}$, the same telescoping argument yields $f(\bar x_T)-f(x^{\star})=O\big(\frac{\log T}{T}\big)$.
\end{itemize}

-----------------------------------------------------------------------

\end{document}